\subsection{Solution Integration}

The final solution is a staged determinant-based workflow rather than a single monolithic training script. Its purpose is to preserve biological traceability across every transformation from raw curated resources to model outputs. At a high level, the integrated workflow begins by obtaining curated genotype determinant records and AST phenotype records, then restricting both resources to the shared isolate cohort. It continues by cleaning phenotype labels and resolving duplicate phenotype rows, cleaning genotype determinant rows and collapsing repeated detections within each isolate, normalizing genotype class/subclass annotations into a single-class, single-subclass representation, and connecting genotype evidence to phenotype under the \textit{strict}, \textit{broad}, and \textit{all} scopes. The workflow then generates one prepared model-input table per antibiotic per scope and evaluates both rule-based and machine-learning predictors on the prepared tables. This integration is important because the final prediction can still be traced back to a sequence of explicit biological decisions, including how phenotype labels were defined, how duplicate genotype records were resolved, how combined class/subclass annotations were normalized, and how much biological scope was allowed before learning. The project therefore treats feature engineering itself as part of the scientific contribution, not merely as implementation detail.
